% ===========================================================================
% Gesamtskript VWL --- Übungsaufgaben Lektion 2
% Wettbewerb, Staatseingriffe & VGR (Rahmenplan DIHK 1.1)
% Alle Situationen, Unternehmen und Zahlen sind frei erfunden.
% ===========================================================================

\subsection*{Aufgabe~2.1 --- UWG oder GWB? (12 Punkte)}
\addcontentsline{toc}{subsection}{Aufgabe~2.1 --- UWG oder GWB?}

Sie sind Mitarbeiterin bzw. Mitarbeiter im Vertrieb der \emph{Rheinquell Getränke GmbH},
einem regionalen Hersteller von Mineralwasser und Erfrischungsgetränken. In der
wöchentlichen Vertriebsrunde werden drei Vorgänge besprochen:

\begin{enumerate}[label=(\arabic*)]
  \item Ein Wettbewerber wirbt auf seinen Etiketten mit dem Hinweis „Testsieger 2026``,
    obwohl es den genannten Test nie gegeben hat.
  \item Ein anderer Wettbewerber lässt Privathaushalte ohne deren vorherige Einwilligung
    anrufen, um ihnen ein Getränke-Abo zu verkaufen.
  \item Der Vertriebsleiter berichtet, dass auf einem Branchentreffen zwei Konkurrenten
    vorgeschlagen haben, die Listenpreise für die Gastronomie zum 1.~Januar gemeinsam
    um 8~\% anzuheben. Er findet die Idee „wirtschaftlich vernünftig``.
\end{enumerate}

\begin{enumerate}[label=\alph*)]
  \item \textbf{Unterscheiden} Sie das Gesetz gegen den unlauteren Wettbewerb (UWG) und
    das Gesetz gegen Wettbewerbsbeschränkungen (GWB) hinsichtlich ihres Schutzzwecks.
    \hfill (4~P.)
  \item \textbf{Ordnen} Sie die Vorgänge (1) bis (3) jeweils dem zutreffenden Gesetz zu und
    \textbf{begründen} Sie Ihre Zuordnung. \hfill (6~P.)
  \item \textbf{Nennen} Sie je zwei mögliche Rechtsfolgen, mit denen Verstöße gegen das UWG
    bzw. gegen das GWB geahndet werden können. \hfill (2~P.)
\end{enumerate}

\subsection*{Aufgabe~2.2 --- Kartellverbot und Ausnahmen (10 Punkte)}
\addcontentsline{toc}{subsection}{Aufgabe~2.2 --- Kartellverbot und Ausnahmen}

Sie sind Assistentin bzw. Assistent der Geschäftsleitung der \emph{Elbtal Tiefbau GmbH}
(85~Beschäftigte) in Sachsen. Die Geschäftsführerin bittet Sie, vier Vorhaben
kartellrechtlich einzuschätzen:

\begin{enumerate}[label=(\arabic*)]
  \item Vor einer Ausschreibung des Landkreises für die Sanierung einer Kreisstraße
    sollen Elbtal und zwei Mitbewerber vereinbaren, dass diesmal Elbtal das günstigste
    Angebot abgibt; bei der nächsten Ausschreibung ist ein anderer Mitbewerber „an der Reihe``.
  \item Elbtal will sich mit vier anderen kleinen Tiefbaubetrieben zu einer
    Einkaufsgemeinschaft für Baumaschinen und Baustoffe zusammenschließen, um gegenüber
    den großen Baukonzernen bessere Einkaufskonditionen zu erzielen. Ihre Preise und
    Angebote gestalten alle Betriebe weiterhin selbst.
  \item Elbtal und ein Wettbewerber wollen gemeinsam ein CO$_2$-reduziertes
    Asphaltverfahren entwickeln. Die Ergebnisse dürfen beide nutzen; ihre Preise legen
    sie getrennt fest. Ihr gemeinsamer Marktanteil liegt bei rund 12~\%.
  \item Ein Wettbewerber aus Brandenburg schlägt vor, dass Elbtal künftig keine Angebote
    mehr in Brandenburg abgibt; im Gegenzug verzichtet er auf Aufträge in Sachsen.
\end{enumerate}

\begin{enumerate}[label=\alph*)]
  \item \textbf{Grenzen} Sie das Kartell von der Fusion hinsichtlich der rechtlichen und
    wirtschaftlichen Selbstständigkeit der beteiligten Unternehmen \textbf{ab}. \hfill (2~P.)
  \item \textbf{Beurteilen} Sie die Vorhaben (1) bis (4) jeweils unter Angabe der Kartellart
    bzw. der einschlägigen Ausnahme, ob sie zulässig sind. \hfill (8~P.)
\end{enumerate}

\subsection*{Aufgabe~2.3 --- Fusionskontrolle und Missbrauchsaufsicht (12 Punkte)}
\addcontentsline{toc}{subsection}{Aufgabe~2.3 --- Fusionskontrolle und Missbrauchsaufsicht}

Sie arbeiten in der Strategieabteilung der \emph{Nordlicht Molkerei AG}. Der Vorstand plant,
sämtliche Anteile der \emph{Hansemilch GmbH} zu übernehmen. Für das letzte Geschäftsjahr
liegen folgende Umsatzerlöse vor:

\begin{center}
\begin{tabular}{l r r}
\toprule
\rowcolor{bfwTabHead} \textbf{Unternehmen} & \textbf{Umsatz weltweit} & \textbf{davon im Inland}\\
\midrule
Nordlicht Molkerei AG & 430~Mio.~€ & 390~Mio.~€\\
Hansemilch GmbH & 85~Mio.~€ & 22~Mio.~€\\
\bottomrule
\end{tabular}
\end{center}

\noindent Auf dem norddeutschen Markt für Frischmilchprodukte hätten beide Unternehmen
nach dem Zusammenschluss einen gemeinsamen Marktanteil von rund 43~\%.

\begin{enumerate}[label=\alph*)]
  \item \textbf{Prüfen} Sie anhand der Umsatzschwellen des \S\,35 GWB, ob der Zusammenschluss
    beim Bundeskartellamt anzumelden ist. \textbf{Prüfen} Sie außerdem, ob sich das Ergebnis
    ändert, wenn die Hansemilch GmbH im Inland nur 15~Mio.~€ umsetzen würde (übrige Werte
    unverändert, keine weiteren Besonderheiten). \hfill (5~P.)
  \item \textbf{Erläutern} Sie, wie das Bundeskartellamt nach der Anmeldung vorgeht und
    unter welcher Voraussetzung es den Zusammenschluss untersagen kann. Gehen Sie dabei auf
    den Marktanteil von 43~\% ein. \hfill (3~P.)
  \item Das Bundeskartellamt untersagt die Fusion. \textbf{Erklären} Sie, welche Möglichkeit
    Nordlicht dennoch hat, den Zusammenschluss zu verwirklichen. \hfill (2~P.)
  \item Angenommen, die Fusion kommt zustande und Nordlicht senkt anschließend den
    Milchabnahmepreis für die von ihr abhängigen Landwirte weit unter das übliche
    Marktniveau --- mit der Drohung, andernfalls keine Milch mehr abzunehmen.
    \textbf{Ordnen} Sie dieses Verhalten der zutreffenden Missbrauchsform \textbf{zu} und
    \textbf{begründen} Sie. \hfill (2~P.)
\end{enumerate}

\subsection*{Aufgabe~2.4 --- Höchstpreis auf dem Pelletmarkt (12 Punkte)}
\addcontentsline{toc}{subsection}{Aufgabe~2.4 --- Höchstpreis auf dem Pelletmarkt}

Sie sind Mitarbeiterin bzw. Mitarbeiter im Controlling der \emph{Pelletwerk Sauerland GmbH}.
Nach stark gestiegenen Energiepreisen diskutiert die Regierung einen gesetzlichen
Höchstpreis für Holzpellets. Für den (vereinfachten) Gesamtmarkt gelten folgende
Funktionen ($p$ = Preis in € je Tonne, $x$ = Menge in 1.000~Tonnen):
\[
\text{Nachfrage: } x_N = 900 - 1{,}5\,p \qquad\qquad \text{Angebot: } x_A = -300 + 2{,}5\,p
\]

\begin{enumerate}[label=\alph*)]
  \item \textbf{Berechnen} Sie Gleichgewichtspreis und Gleichgewichtsmenge. \hfill (3~P.)
  \item Der Höchstpreis wird auf 240~€/t festgelegt. \textbf{Berechnen} Sie die angebotene
    und die nachgefragte Menge und \textbf{bestimmen} Sie Art und Umfang der entstehenden
    Marktsituation. \textbf{Begründen} Sie außerdem, ob ein Höchstpreis von 320~€/t eine
    Wirkung hätte. \hfill (4~P.)
  \item Es bildet sich ein Schwarzmarkt. \textbf{Berechnen} Sie den Preis, den die
    Nachfrager für die tatsächlich angebotene Menge höchstens zu zahlen bereit wären, und
    \textbf{erläutern} Sie zwei Maßnahmen, die der Staat flankierend ergreifen müsste,
    damit der Höchstpreis seine Wirkung entfalten kann. \hfill (3~P.)
  \item \textbf{Berechnen} Sie, wie sich der Gesamtumsatz aller Anbieter durch den
    Höchstpreis absolut und prozentual verändert. \hfill (2~P.)
\end{enumerate}

\subsection*{Aufgabe~2.5 --- Marktkonforme und marktkonträre Eingriffe (12 Punkte)}
\addcontentsline{toc}{subsection}{Aufgabe~2.5 --- Marktkonforme und marktkonträre Eingriffe}

Sie arbeiten in der Stabsstelle Energiewirtschaft der \emph{Stadtwerke Lahntal GmbH}. Um
private Haushalte bei den Stromkosten zu entlasten, werden politisch vier Vorschläge
diskutiert:

\begin{enumerate}[label=(\arabic*)]
  \item eine gesetzliche Obergrenze von 30~Cent je kWh für Haushaltsstrom, die kein
    Versorger überschreiten darf --- ohne finanziellen Ausgleich für die Versorger,
  \item eine Senkung der Stromsteuer auf das EU-rechtlich zulässige Mindestmaß,
  \item ein einmaliger Energiekostenzuschuss für Haushalte mit geringem Einkommen,
  \item eine staatliche Investitionszulage für Hersteller von Wärmepumpen.
\end{enumerate}

\begin{enumerate}[label=\alph*)]
  \item \textbf{Ordnen} Sie die Vorschläge (1) bis (4) den marktkonformen bzw.
    marktkonträren Eingriffen \textbf{zu} und \textbf{begründen} Sie jeweils kurz.
    \hfill (4~P.)
  \item \textbf{Erläutern} Sie die Wirkung von Vorschlag (2) auf Angebot,
    Gleichgewichtspreis und Gleichgewichtsmenge am Strommarkt. \hfill (3~P.)
  \item Die Obergrenze aus Vorschlag (1) liegt unter dem Gleichgewichtspreis.
    \textbf{Erläutern} Sie drei mögliche Folgen für die Stadtwerke Lahntal bzw. den
    Strommarkt. \hfill (3~P.)
  \item \textbf{Beurteilen} Sie aus ordnungspolitischer Sicht, ob Vorschlag (3) oder
    Vorschlag (1) vorzuziehen ist. \hfill (2~P.)
\end{enumerate}

\subsection*{Aufgabe~2.6 --- VGR: Vom BIP zum Volkseinkommen (16 Punkte)}
\addcontentsline{toc}{subsection}{Aufgabe~2.6 --- VGR: Vom BIP zum Volkseinkommen}

Sie sind Mitarbeiterin bzw. Mitarbeiter der Research-Abteilung der \emph{Nordbank Consult
GmbH} und bereiten für einen Kundenvortrag die Daten der (fiktiven) Volkswirtschaft
\emph{Aldoria} auf. Für das abgelaufene Jahr liegen folgende Werte vor (in Mrd.~€):

\begin{center}
\begin{tabular}{l r}
\toprule
\rowcolor{bfwTabHead} \textbf{Größe} & \textbf{Mrd.~€}\\
\midrule
Private Konsumausgaben & 2.180\\
Konsumausgaben des Staates & 860\\
Bruttoinvestitionen & 740\\
Exporte von Waren und Dienstleistungen & 1.950\\
Importe von Waren und Dienstleistungen & 1.730\\
Primäreinkommen, die Inländer aus der übrigen Welt erhalten & 210\\
Primäreinkommen, die an die übrige Welt (Nichtinländer) fließen & 145\\
Abschreibungen & 690\\
Produktions- und Importabgaben & 420\\
Subventionen & 70\\
Arbeitnehmerentgelt & 1.936\\
\bottomrule
\end{tabular}
\end{center}

\begin{enumerate}[label=\alph*)]
  \item \textbf{Berechnen} Sie den Außenbeitrag und das Bruttoinlandsprodukt (BIP) über
    die Verwendungsrechnung. \hfill (4~P.)
  \item \textbf{Berechnen} Sie das Bruttonationaleinkommen (BNE) und \textbf{erläutern} Sie,
    warum es in Aldoria größer ist als das BIP. \hfill (3~P.)
  \item \textbf{Berechnen} Sie das Nettonationaleinkommen und das Volkseinkommen.
    \hfill (4~P.)
  \item \textbf{Berechnen} Sie die Unternehmens- und Vermögenseinkommen sowie die Lohn- und
    die Gewinnquote. \hfill (3~P.)
  \item Im Folgejahr steigt die Lohnquote auf 66~\%, obwohl die Arbeitnehmerentgelte
    unverändert geblieben sind. \textbf{Erläutern} Sie, wie dies zu erklären ist.
    \hfill (2~P.)
\end{enumerate}

\subsection*{Aufgabe~2.7 --- Nominales und reales Wachstum (12 Punkte)}
\addcontentsline{toc}{subsection}{Aufgabe~2.7 --- Nominales und reales Wachstum}

Sie sind im Controlling der \emph{Weserland Logistik GmbH} tätig. Die Geschäftsführung
wundert sich: „Das BIP steigt jedes Jahr --- warum sprechen die Medien trotzdem von einer
schrumpfenden Wirtschaft?`` Sie erhalten folgende Daten (Jahr~1 = Basisjahr):

\begin{center}
\begin{tabular}{l r r r}
\toprule
\rowcolor{bfwTabHead} & \textbf{Jahr 1} & \textbf{Jahr 2} & \textbf{Jahr 3}\\
\midrule
Nominales BIP (Mrd.~€) & 4.000,0 & 4.202,4 & 4.263,0\\
BIP-Deflator (Basisjahr = 100) & 100,0 & 103,0 & 105,0\\
\bottomrule
\end{tabular}
\end{center}

\begin{enumerate}[label=\alph*)]
  \item \textbf{Berechnen} Sie das reale BIP für die Jahre~2 und~3. \hfill (3~P.)
  \item \textbf{Berechnen} Sie die nominalen und die realen Wachstumsraten für die
    Jahre~2 und~3 (auf zwei Nachkommastellen). \hfill (4~P.)
  \item \textbf{Berechnen} Sie die Preissteigerung (gemessen am Deflator) von Jahr~2 auf
    Jahr~3 und \textbf{erläutern} Sie damit die Frage der Geschäftsführung. \hfill (3~P.)
  \item \textbf{Unterscheiden} Sie den BIP-Deflator vom Verbraucherpreisindex (VPI).
    \hfill (2~P.)
\end{enumerate}

\subsection*{Aufgabe~2.8 --- Magisches Viereck: Lage beurteilen (14 Punkte)}
\addcontentsline{toc}{subsection}{Aufgabe~2.8 --- Magisches Viereck: Lage beurteilen}

Sie sind in der Geschäftsleitungsassistenz der \emph{Allgäu Präzisionstechnik GmbH}, eines
stark exportorientierten Maschinenbauers, tätig. Für die Strategieklausur sollen Sie die
gesamtwirtschaftliche Lage eines Landes beurteilen, in dem das Unternehmen einen großen
Teil seines Umsatzes erzielt. Es liegen folgende Jahreswerte vor:

\begin{center}
\begin{tabular}{l r}
\toprule
\rowcolor{bfwTabHead} \textbf{Kennzahl} & \textbf{Wert}\\
\midrule
Inflationsrate (HVPI) & 4,1~\%\\
Arbeitslosenquote & 3,4~\%\\
Wachstumsrate des realen BIP & 0,2~\%\\
Außenbeitrag in \% des nominalen BIP & 5,5~\%\\
\bottomrule
\end{tabular}
\end{center}

\begin{enumerate}[label=\alph*)]
  \item \textbf{Beurteilen} Sie für jedes der vier Ziele des magischen Vierecks anhand der
    Kennzahlen, ob es erreicht ist. Nennen Sie dabei jeweils den Zielwert. \hfill (8~P.)
  \item Die Regierung will die Inflation durch deutliche Kürzungen der Staatsausgaben
    bekämpfen. \textbf{Erläutern} Sie, welcher Zielkonflikt dadurch entstehen kann.
    \hfill (3~P.)
  \item Der Geschäftsführer meint: „Ein Außenbeitrag von 5,5~\% zeigt doch nur, wie stark
    das Land ist --- je höher, desto besser.`` \textbf{Beurteilen} Sie diese Aussage.
    \hfill (3~P.)
\end{enumerate}
